\documentclass[11pt,a4paper]{article}

\usepackage[utf8]{inputenc}
\usepackage[T1]{fontenc}
\usepackage{amsmath,amssymb}
\usepackage{geometry}
\usepackage{booktabs}
\usepackage{listings}
\usepackage[colorlinks=true,linkcolor=blue,citecolor=blue,urlcolor=blue]{hyperref}

\geometry{left=2.5cm,right=2.5cm,top=2.5cm,bottom=2.5cm}

\lstset{
  basicstyle=\ttfamily\small,
  breaklines=true,
  frame=single,
  numbers=left,
  numberstyle=\tiny,
  tabsize=2
}

\title{\textbf{NRR-Coupled (cp): Coupled State Update Under Candidate Dependencies}}
\author{Kei Saito\\Independent Researcher, Japan\\\texttt{kei.saito.research@gmail.com}}
\date{February 2026\\[0.5em]\textit{Part of the Non-Resolution Reasoning (NRR) research program.}}

\begin{document}
\maketitle

\begin{center}
\textcopyright\ 2026 Kei Saito. Licensed under CC BY 4.0.\\
\url{https://creativecommons.org/licenses/by/4.0/}
\end{center}

\begin{abstract}
NRR-Transfer established an uncoupled state-update interface for independent candidates: the LLM selects a directional operator (\(\sigma\)/\(\delta\)) and the client executes deterministic weight updates. This paper defines NRR-Coupled (cp), a coupled extension for dependent-candidate conditions while keeping the same LLM decision interface. The new component is a client-side propagation rule driven by a dependency matrix. We provide a formal update rule, per-turn complexity bounds, applicability conditions, and explicit failure conditions. The main claim is conditional: under dependency-present conditions, coupled propagation can improve state-tracking utility relative to the uncoupled baseline; under independent conditions, coupled propagation is not expected to improve and may regress if misapplied. The paper includes a falsifiable evaluation contract and a reproducibility procedure.
\end{abstract}

\noindent\textbf{Keywords:} NRR-Coupled, coupled update, dependency propagation, deterministic client execution, falsifiable evaluation

\section{Introduction}
LLM APIs are stateless across calls, but many systems require state updates over candidate sets. NRR-Transfer addressed this with an uncoupled interface under independent-candidate assumptions. This paper introduces \textbf{NRR-Coupled (cp)} for conditions where candidates are not independent.

NRR is not an anti-LLM framework.
NRR does not replace standard LLM use.
NRR optimizes when to commit and when to defer, under explicit conditions.

\textbf{Claim boundary.}
Transfer covers independent-candidate conditions; cp covers dependent-candidate conditions. cp does not redefine the uncoupled baseline and treats the Transfer definition as canonical.

\textbf{Design constraint (first hypothesis).}
The LLM-side selection of \(\sigma\)/\(\delta\) and target index is unchanged. Only the client propagation rule changes. If this interface assumption fails, we record a re-evaluation flag for Transfer/Principles and continue cp design.

Series Consistency Statement: Unsupported interpretations are never injected at a fixed evidence state. Hypothesis-set expansion is evidence-gated, and unobserved possibilities remain unresolved. Evidence in this paper is restricted to explicit input, provided context, and declared dependency structure.

\section{Scope and Notation}
Let \(V=\{1,\dots,n\}\) be candidates and \(\mathbf{w}_t\in\mathbb{R}^n\) be weights at turn \(t\), with:
\[
\sum_{i=1}^n w_t^{(i)} = 1,\qquad w_t^{(i)}\ge \varepsilon > 0.
\]

Dependency structure is represented by matrix \(A\in\mathbb{R}^{n\times n}\), where \(A_{ij}\) denotes influence from candidate \(j\) to candidate \(i\):
\begin{itemize}
  \item \(A_{ij}>0\): reinforcing relation,
  \item \(A_{ij}<0\): competitive relation,
  \item \(A_{ij}=0\): no direct dependency.
\end{itemize}

\textbf{Fixed interface contract.}
At each turn, the LLM emits \((o_t,k_t)\), where \(o_t\in\{\sigma,\delta\}\) and \(k_t\in V\). Client execution is deterministic.

\section{Coupled Update Rule (cp New Definition)}
\subsection{Step 1: Local update (Transfer-compatible)}
Define step size \(\alpha>0\), clipping upper bound \(u_{\max}\le 1-\varepsilon\), and:
\[
\Delta_t = \begin{cases}
+\alpha & (o_t=\sigma),\\
-\alpha & (o_t=\delta).
\end{cases}
\]

Local provisional update:
\[
\tilde{w}^{(k_t)} = \operatorname{clip}\left(w_t^{(k_t)}+\Delta_t,\,\varepsilon,\,u_{\max}\right),\qquad
\tilde{w}^{(i)} = w_t^{(i)}\;(i\neq k_t).
\]

\subsection{Step 2: Coupled propagation (cp extension)}
Let \(d_t=\tilde{w}^{(k_t)}-w_t^{(k_t)}\) and coupling gain \(\beta\ge 0\). For \(i\neq k_t\):
\[
\tilde{w}^{(i)} \leftarrow \operatorname{clip}\left(\tilde{w}^{(i)} + \beta A_{i,k_t} d_t,\,\varepsilon,\,u_{\max}\right).
\]

Equivalent vector form (with sparse execution):
\[
\tilde{\mathbf{w}} \leftarrow \operatorname{clip}\left(\tilde{\mathbf{w}} + \beta\,A_{:,k_t}d_t,\,\varepsilon,\,u_{\max}\right),
\]
with the target index excluded from propagation overwrite.

\subsection{Step 3: Renormalization}
\[
S_t = \sum_{i=1}^{n}\tilde{w}^{(i)},\qquad
w_{t+1}^{(i)} = \frac{\tilde{w}^{(i)}}{S_t}.
\]
Since clipping enforces positivity, \(S_t>0\) always holds.

\subsection{Pseudocode}
\begin{lstlisting}[caption={cp single-turn client update (fixed LLM interface)}]
Inputs:
  w[1..n], A[1..n,1..n], operator o in {sigma, delta}, target k
  alpha > 0, beta >= 0, epsilon, u_max

# Step 1: local update (same interface as Transfer)
if o == sigma:
    delta = +alpha
else:
    delta = -alpha

old_k = w[k]
w[k] = clip(w[k] + delta, epsilon, u_max)
d = w[k] - old_k

# Step 2: coupled propagation (new in cp)
for i in neighbors_from_column_k(A, k):
    if i == k:
        continue
    w[i] = clip(w[i] + beta * A[i, k] * d, epsilon, u_max)

# Step 3: renormalization
S = sum(w)
for i in 1..n:
    w[i] = w[i] / S

return w
\end{lstlisting}

\section{Complexity Upper Bounds}
Let \(\deg_{\text{out}}(k_t)\) be the number of nonzero entries in column \(k_t\) of \(A\).

Per turn complexity:
\begin{itemize}
  \item local update: \(O(1)\)
  \item coupled propagation (sparse): \(O(\deg_{\text{out}}(k_t))\)
  \item renormalization: \(O(n)\)
\end{itemize}

Therefore,
\[
T_{\text{turn}} = O\big(n + \deg_{\text{out}}(k_t)\big).
\]
Dense-column worst case gives \(\deg_{\text{out}}(k_t)=n-1\), so optimized execution is \(O(n)\) per turn.
A naive full-matrix multiply implementation is bounded by \(O(n^2)\) per turn.

Memory complexity:
\[
M = O(n+\operatorname{nnz}(A))\quad\text{(sparse)}\qquad\text{or}\qquad O(n^2)\quad\text{(dense)}.
\]

\section{Applicability and Failure Conditions}
\subsection{Applicability (required conditions)}
cp is applicable under tested conditions when all of the following hold:
\begin{enumerate}
  \item finite trackable candidate set,
  \item dependency structure can be specified or estimated as bounded matrix \(A\),
  \item directional LLM decision (\(\sigma/\delta\), target) remains sufficient,
  \item coupling gain \(\beta\) is bounded to keep updates stable in observed runs.
\end{enumerate}

\subsection{Failure conditions (must be reported)}
\begin{enumerate}
  \item \textbf{Oscillation}: repeated sign-flip drift in key candidates over consecutive turns.
  \item \textbf{Over-coupling collapse}: one candidate dominates due to excessive \(\beta\) or wrong \(A\).
  \item \textbf{Graph mismatch}: cp underperforms uncoupled baseline in dependency-present tasks.
  \item \textbf{Interface break}: parser cannot map LLM output to valid \((\sigma/\delta, k)\).
\end{enumerate}

If (3) or (4) occurs, set `TRANSFER\_PRINCIPLES\_REEVAL\_REQUIRED=true` and continue cp-side diagnosis rather than stopping the design process.

\section{Evaluation Plan (Falsifiable)}
\subsection{Paired comparison}
All runs compare two systems under identical prompts, models, temperature, and seeds:
\begin{itemize}
  \item \textbf{Uncoupled baseline}: Transfer reference definition.
  \item \textbf{Coupled}: cp rule introduced in this paper.
\end{itemize}

\subsection{Datasets/conditions}
\begin{itemize}
  \item \textbf{D-dependent}: scenarios with known nonzero dependencies.
  \item \textbf{D-independent}: scenarios where candidates are effectively independent.
  \item \textbf{D-mismatch}: intentionally perturbed dependency graph.
\end{itemize}

\subsection{Primary metrics}
\begin{itemize}
  \item utility score \(U\): task-defined objective normalized to \([0,1]\),
  \item correction-turn rate \(C\): fraction of turns spent undoing prior state drift,
  \item instability rate \(R_{\text{osc}}\): fraction of runs triggering oscillation criterion,
  \item token cost per turn \(K\).
\end{itemize}

\subsection{Pre-registered thresholds (fixed before results)}
\begin{center}
\begin{tabular}{@{}ll@{}}
\toprule
Condition & Criterion \\
\midrule
Success-1 (D-dependent) & $U_{cp} - U_{base} \ge 0.05$ \\
Success-2 (D-independent) & $U_{cp} - U_{base} \ge -0.02$ \\
Success-3 (stability) & $R_{osc,cp} \le 0.01$ \\
Failure-A & Success-1 not met \\
Failure-B & Success-2 violated \\
Failure-C & Success-3 violated \\
\bottomrule
\end{tabular}
\end{center}

These criteria are falsifiable and can reject cp under the tested setup.

\section{Reproducibility Protocol}
The repository includes:
\begin{itemize}
  \item specification: \texttt{spec/nrr-coupled\_spec\_v1.md}
  \item simulation script: \texttt{repro/coupled\_state\_sim.py}
  \item execution steps: \texttt{repro/README.md}
\end{itemize}

No post-hoc threshold tuning is allowed after observing outcomes.

\section{Limitations}
This version is a specification and evaluation-plan manuscript. It does not claim universal superiority of coupled updates. Gains are expected only under dependency-present conditions and may disappear or reverse under independence or graph misspecification.

\section{Conclusion}
cp extends Transfer by adding a deterministic client-side coupled propagation rule while keeping the same LLM interface. The contribution is bounded: a formal update law, complexity bounds, explicit applicability/failure conditions, and a falsifiable evaluation contract for dependency-aware state management.

\begin{thebibliography}{9}
\bibitem{nrr_transfer_2026}
K. Saito,
\textit{NRR-Transfer: Cross-Domain Transfer of Phase 1.5 Operators Under Fixed Interface Constraints},
2026.

\bibitem{nrr_phi_2026}
K. Saito,
\textit{NRR-Phi: Text-to-State Mapping and Operator Principles},
2026.

\bibitem{nrr_core_2025}
K. Saito,
\textit{Non-Resolution Reasoning (NRR): A Framework for Conditional Commitment},
2025.
\end{thebibliography}

\end{document}
